\chapter{What the existing software already supplies}
\label{gap29:software-followup}
The public experiment demonstrates why reuse is the starting point. Native
R fitting and a separate Python calculation agree on the same observations,
regression specification and HC2 covariance. Rewriting a familiar estimator
in another language would add little to that calculation. The proposal's
contribution must instead be demonstrated in the connection between the
economic question, assumptions, available information and reported decision.

\section{Panel estimation is already a commercial capability}
The official MATLAB documentation retrieved on 29 October 2026 describes
one-way fixed- and random-effects panel regressions. The fixed-effects
documentation explains transformed-data estimation, a subject-clustered
sandwich option and handling of missing measurements. It documents
\texttt{fitfepanel} as introduced in R2026b and
\texttt{fitrepanel} in R2026a. The retrieved pages are current R2026b pages;
their exact availability before the 7 October literature cutoff has not been
established. They are a dated capability update, not a reconstructed
pre-cutoff software comparison. \citep{gap29matlab}

The documented treatment of missing identifiers is consequential: a record
with an absent subject identifier is assigned an additional group, while a
missing predictor or response removes the observation from estimation.
Those are computational conventions, not evidence that the economic unit
is known or that discarded observations are ignorable. The panel study in
Chapter~\ref{gap29:funding} shows why that distinction matters.

A matched comparison must specify the institution and time effects,
covariance convention, initial sample and missing-record treatment before
comparing coefficients. The inspected pages describe a one-way model and
subject clustering; they are not sufficient evidence about every possible
way to obtain time effects or two-way covariance in the wider MATLAB estate.
There has been no licensed MATLAB execution in this study. The appropriate
proposal decision is to retain it as an existing-estate alternative and
test the declared calculation where the institution already has access.
Claims of missing commercial capability would require broader evidence.

\section{A decision diagram already contains an information structure}
HUGIN's official GUI 9.8 manual describes Bayesian networks, decision and
utility nodes, and limited-memory influence diagrams. Its account makes
the information available at each decision explicit through incoming
informational links; it also relaxes the traditional requirement that every
earlier observation remains available. This is directly relevant to a
servicing policy, where later actions must use the information actually
observed at their dates. \citep{gap29hugin}

The proposal should therefore compare a graph assistant with an established
diagram editor and a declared decision model. It should not equate drawing
an arrow with identifying an intervention effect, or infer from a manual
that the proposed integration has been implemented. The inspected material
supports a concrete alternative for representing decisions and information;
it does not establish an automatic solution to latent confounding or missing
outcomes.

The vendor's published licence text makes distribution using its API subject
to an OEM licence. That term changes the proposed reuse route: an internal
connection to an entitled installation and redistribution of a new product
are different propositions. The actual order, deployment conditions and
prices remain to be established with the vendor. The current manual's exact
release date relative to the literature cutoff is also unresolved.
\citep{gap29huginlicense}

\section{Which comparisons belong before a prototype}
The proposal can justify selecting existing estimators and diagram tools
from inspected methods, alternatives and explicit limitations. A complete
installed distribution, matched adverse-data tests and target-system
acceptance are subsequent implementation evidence. They must be specified
before a native adapter is accepted, but they should not be treated as
proof that the economic case for a prototype is unfinished in every respect.

DoubleML, \texttt{did}, \texttt{causaleffect}, and the named commercial
causaLens offering still require their previously identified version,
method or reuse questions to be resolved where those questions could change
the selection. The present calculations neither execute those packages
nor certify their adapters. The commercial source follow-up narrows two
documentation gaps; it does not complete the 102-cell software census.
The decision remains to compare a narrow integration with the institution's
existing practice, and to spend on implementation only where that comparison
can reveal a material analytical benefit.
